%%%%%%%%%%%%%%%%%%%%%%%%%%%%%%%%%%%%%%%%%%%%%%%%%%%%%%%
% NOTAS TEORIA DE TRANSPORTE OPTIMO
%
% Contains ?? figures. 
% File started: ??
% First circulated version: ??
% Last version:
% Last modification: -
%%%%%%%%%%%%%%%%%%%%%%%%%%%%%%%%%%%%%%%%%%%%%%%%%%%%%%%

\documentclass[11pt,a4paper]{amsart}
\usepackage{xr-hyper}
\newcounter{chapter}

\usepackage{amssymb}
\usepackage{hyperref}
\usepackage{listings}
\usepackage{mathtools}
\usepackage{enumitem}

% Fuzz -------------------------------------------------------------------
\hfuzz2pt % Don't bother to report over-full boxes if over-edge is < 2pt
% Line spacing -----------------------------------------------------------

%%%%%%%%%fin pegada%%%%%%%%%%%%%%%%%%5

\usepackage[T1]{fontenc}
\usepackage[spanish]{babel}
\usepackage{graphicx}
\usepackage{tikz}
\usetikzlibrary{babel}

\setlength{\textwidth}{15.5cm} 
\setlength{\textheight}{24cm}
\setlength{\oddsidemargin}{0cm}
\setlength{\evensidemargin}{0cm}


\parskip 3pt

%%
%% OPERADORES
%%
\renewcommand{\Re}{\operatorname{Re}}
\newcommand{\tr}{\operatorname{tr}}
\newcommand{\id}{\operatorname{id}}
\def\div{\operatorname {\mathrm{div}}}
\def\Dom{\operatorname {\mathrm{Dom}}}
\def\argmin{\operatorname {\mathrm{arg\, min}}}
\def\argmax{\operatorname {\mathrm{arg\, max}}}
\def\gen{\operatorname {\mathrm{gen}}}
\def\Im{\operatorname {\mathrm{Im}}}
\def\Nu{\operatorname {\mathrm{Nu}}}
\def\sop{\operatorname {\mathrm{sop}}}
\def\diam{\operatorname {\mathrm{diam}}}
\def\dist{\operatorname {\mathrm{dist}}}
\def\var{\operatorname {\mathrm{Var}}}
\def\codim{\operatorname {\mathrm{codim}}}
\def\pint{\operatorname {--\!\!\!\!\!\int\!\!\!\!\!--}}

%%
%% EPSILON
%%
\newcommand{\ve}{\varepsilon}

%%
%% LETRAS MATHBB
%%
\newcommand{\R}{{\mathbb R}}
\newcommand{\Z}{{\mathbb Z}}
\newcommand{\C}{{\mathbb C}}
\newcommand{\Q}{{\mathbb Q}}
\newcommand{\N}{{\mathbb N}}
\newcommand{\E}{{\mathbb E}}
\newcommand{\V}{{\mathbb V}}
\renewcommand{\P}{{\mathbb P}}
\newcommand{\II}{{\mathbb I}}
\newcommand{\OO}{{\mathbb O}}
\newcommand{\Sb}{{\mathbb S}}
\newcommand{\Tb}{{\mathbb T}}

%%
%% LETRAS CALIFGRAFICAS
%%
\newcommand{\F}{{\mathcal F}}
\newcommand{\Hc}{{\mathcal H}}
\newcommand{\Tc}{{\mathcal T}}
\newcommand{\Ec}{{\mathcal E}}
\newcommand{\Vc}{{\mathcal V}}
\newcommand{\A}{{\mathcal A}}
\newcommand{\J}{{\mathcal J}}
\newcommand{\G}{{\mathcal G}}
\newcommand{\U}{{\mathcal U}}
\newcommand{\B}{{\mathcal B}}
\newcommand{\Sw}{{\mathcal S}}
\newcommand{\D}{{\mathcal D}}
\newcommand{\M}{{\mathcal M}}
\newcommand{\LL}{{\mathcal L}}

%%
%% LETRAS MATHBF
%%
\newcommand{\ab}{{\mathbf a}}
\newcommand{\bb}{{\mathbf b}}
\newcommand{\cc}{{\mathbf c}}
\newcommand{\ee}{{\mathbf e}}
\newcommand{\ff}{{\mathbf f}}
\newcommand{\gb}{{\mathbf g}}
\newcommand{\hh}{{\mathbf h}}
\newcommand{\n}{{\mathbf n}}
\newcommand{\mm}{{\mathbf m}}
\newcommand{\pp}{{\mathbf p}}
\newcommand{\qq}{{\mathbf q}}
\newcommand{\sbf}{{\mathbf s}}
\newcommand{\tb}{{\mathbf t}}
\newcommand{\uu}{{\mathbf u}}
\newcommand{\vv}{{\mathbf v}}
\newcommand{\ww}{{\mathbf w}}
\newcommand{\xx}{{\mathbf x}}
\newcommand{\yy}{{\mathbf y}}
\newcommand{\zz}{{\mathbf z}}
\newcommand{\I}{{\mathbf 1}}
\newcommand{\Pb}{{\mathbf P}}
\newcommand{\XX}{{\mathbf X}}
\newcommand{\YY}{{\mathbf Y}}

%%
%% LETRAS MATHFRAK
%%

\newcommand{\Gk}{{\mathfrak G}}


%%
%% TEOREMAS
%%
\newtheorem{teo}{Teorema}[section]
\newtheorem{lema}[teo]{Lema}
\newtheorem{prop}[teo]{Proposición}
\newtheorem{cor}[teo]{Corolario}

%%
%% REMARK
%%
\theoremstyle{remark}
\newtheorem{obs}[teo]{Observación}

%%
%% DEFINICION
%%
\theoremstyle{definition}
\newtheorem{defi}[teo]{Definición}
\newtheorem{ejem}[teo]{Ejemplo}
\newtheorem{ejer}[teo]{Ejercicio}
\newtheorem{nota}[teo]{Notación}
\newtheorem{ejercicio}{Ejercicio}[chapter]


%%
%% NUMERACION
%%
\numberwithin{subsection}{section}
\numberwithin{equation}{section}

%%
%% TITULO
%%


\externaldocument[N-]{../notas/Notas-TO}
\newcommand{\cuadernolink}[2]{\noindent\textbf{Cuaderno:} \emph{#1} (\href{https://colab.research.google.com/github/jfbonder/transporte-optimo/blob/main/cuadernos/#2}{abrir en Colab}).\par\medskip}
\pagestyle{plain}

\begin{document}
\renewcommand{\thechapter}{7}
\begin{center}
{\large\sc Teoría de Transporte Óptimo}\\[2pt]
{\small 2do cuatrimestre 2026 --- Julián Fernández Bonder}\\[10pt]
{\LARGE\bf Práctica 7}\\[4pt]
{\large Algunas interpretaciones y extensiones}
\end{center}
\bigskip
\noindent{\small Los ejercicios son los de la sección de ejercicios del Capítulo~7 de las notas del curso, con la misma numeración. Las referencias a teoremas, proposiciones y secciones remiten a las notas (\url{https://github.com/jfbonder/transporte-optimo}).}
\medskip


\begin{ejercicio}[Acoplamientos y covarianza]\label{ejer.acoplamiento.covarianza}
Sean $\mu,\nu\in\mathcal P_2(\R)$ con medias $m_\mu,m_\nu$ y varianzas
$\sigma_\mu^2,\sigma_\nu^2$, y sea $(X,Y)$ un vector aleatorio con ley
$\pi\in\Pi(\mu,\nu)$.
\begin{enumerate}[label=(\alph*)]
\item Probar que
$\E|X-Y|^2=(m_\mu-m_\nu)^2+\sigma_\mu^2+\sigma_\nu^2-2\operatorname{Cov}(X,Y)$,
de modo que minimizar el costo cuadrático equivale a maximizar la
covarianza (o la correlación) del acoplamiento.
\item Probar la cota de Cauchy--Schwarz
$\operatorname{Cov}(X,Y)\le\sigma_\mu\sigma_\nu$ y deducir que
\[
\min_{\pi\in\Pi(\mu,\nu)}\E|X-Y|^2\ge(m_\mu-m_\nu)^2+(\sigma_\mu-\sigma_\nu)^2 .
\]
\item Probar que hay igualdad en (b) si y sólo si $Y=aX+b$ casi seguramente
con $a>0$; por lo tanto la cota de (b) se alcanza si y sólo si $\nu$ es la
imagen de $\mu$ por una afinidad creciente (por ejemplo, si ambas son
gaussianas). En el Capítulo~\ref{N-cap.1D} se identifica el acoplamiento que
realiza el mínimo en general.
\end{enumerate}
\end{ejercicio}

\begin{ejercicio}[Cotas de Fréchet--Hoeffding]\label{ejer.frechet.hoeffding}
Sean $\mu,\nu\in\mathcal P(\R)$ con funciones de distribución $F_\mu$, $F_\nu$,
y $\pi\in\Pi(\mu,\nu)$ con función de distribución conjunta
$H(x,y)=\pi((-\infty,x]\times(-\infty,y])$.
\begin{enumerate}[label=(\alph*)]
\item Probar que
$\max\{F_\mu(x)+F_\nu(y)-1,\,0\}\le H(x,y)\le\min\{F_\mu(x),F_\nu(y)\}$
para todo $(x,y)$.
\item Probar que ambas cotas son funciones de distribución de elementos de
$\Pi(\mu,\nu)$. (La cota superior es el plan comonótono del
Capítulo~\ref{N-cap.1D}; la inferior, el ``antimonótono''.)
\end{enumerate}
\end{ejercicio}

\begin{ejercicio}[Matching estable: un ejemplo]\label{ejer.matching.ejemplo}
Dos trabajadores $\{x_1,x_2\}$ y dos empresas $\{y_1,y_2\}$ generan, al
emparejarse, el excedente $U(x_i,y_j)=u_{ij}$ con
$u=\bigl(\begin{smallmatrix}4&3\\3&1\end{smallmatrix}\bigr)$. Todos los
agentes tienen el mismo peso.
\begin{enumerate}[label=(\alph*)]
\item Hallar el matching que maximiza el excedente total y compararlo con
el matching ``codicioso'' que asigna primero el par de mayor excedente.
\item Hallar un sistema de precios $(\varphi,\psi)$ (salarios y beneficios)
con $\varphi_i+\psi_j\ge u_{ij}$ para todo par y con igualdad en los pares
emparejados. Verificar que es un equilibrio: ningún par no emparejado puede
repartirse su excedente de modo que ambos mejoren.
\item Describir todos los sistemas de precios de equilibrio y observar que
forman un segmento, aunque el matching óptimo sea único. ¿Quién se apropia
del excedente en cada extremo del segmento?
\end{enumerate}
\end{ejercicio}

\begin{ejercicio}[Un criterio de optimalidad multimarginal]\label{ejer.multimarginal.criterio}
Sean $\mu_1,\dots,\mu_N\in\mathcal P(\R^d)$ con soporte compacto y
$\gamma\in\Pi(\mu_1,\dots,\mu_N)$ un plan multimarginal concentrado en el
conjunto $\{(x_1,\dots,x_N):\ x_1+\dots+x_N=v\}$ para algún $v\in\R^d$.
\begin{enumerate}[label=(\alph*)]
\item Probar que necesariamente $v=\sum_i\int x\,d\mu_i$.
\item Probar que $\gamma$ es óptimo para todo costo de la forma
$c(x_1,\dots,x_N)=h(x_1+\dots+x_N)$ con $h\colon\R^d\to\R$ convexa.
\emph{Sugerencia:} desigualdad de Jensen.
\item Probar que $\gamma$ es óptimo para el costo
$c(x_1,\dots,x_N)=-\sum_{i\ne j}|x_i-x_j|^2$. \emph{Sugerencia:} expandir
los cuadrados y observar que $\sum_i|x_i|^2$ tiene integral fija sobre
$\Pi(\mu_1,\dots,\mu_N)$, de modo que el problema equivale a (b) con
$h(z)=|z|^2$.
\item ¿Existe siempre un plan como $\gamma$? Estudiar el caso $N=2$,
$\mu_1=\mu_2=\frac12(\delta_0+\delta_1)$, y el caso $N=2$ con $\mu_1$
uniforme en $[0,1]$ y $\mu_2=\delta_0$.
\end{enumerate}
\end{ejercicio}

\begin{ejercicio}[Discretización y estabilidad]\label{ejer.discretizacion.estabilidad}
Sean $X$, $Y$ compactos, $c\in C(X\times Y)$, y $\mu_n\rightharpoonup\mu$,
$\nu_n\rightharpoonup\nu$.
\begin{enumerate}[label=(\alph*)]
\item Dado $\gamma\in\Pi(\mu,\nu)$, probar que existen
$\gamma_n\in\Pi(\mu_n,\nu_n)$ con $\gamma_n\rightharpoonup\gamma$.
\emph{Sugerencia:} sean $\alpha_n\in\Pi(\mu_n,\mu)$ y $\beta_n\in\Pi(\nu,\nu_n)$
planes óptimos para el costo $d(x,x')^2$ (o cualquier costo continuo que
se anule exactamente en la diagonal); pegar $\alpha_n$, $\gamma$ y $\beta_n$
con el Lema~\ref{N-lema.pegado} y proyectar sobre las coordenadas
extremas. Probar que $\int d(x,x')^2\,d\alpha_n\to0$ y usarlo para pasar al límite.
\item Deducir que $\mathcal T_c(\mu_n,\nu_n)\to\mathcal T_c(\mu,\nu)$, donde
$\mathcal T_c$ denota el costo óptimo, y más generalmente que
$\mathcal T_{c_n}(\mu_n,\nu_n)\to\mathcal T_c(\mu,\nu)$ si $c_n\to c$
uniformemente.
\item Concluir que si $\mu_n$, $\nu_n$ son discretizaciones (por ejemplo,
medidas empíricas) de $\mu$, $\nu$, el costo óptimo discreto calculado por
programación lineal converge al costo óptimo continuo, y que los planes
discretos óptimos convergen, salvo subsucesión, a planes óptimos continuos.
\end{enumerate}
\end{ejercicio}

\end{document}
